\documentclass[11pt]{article}
\usepackage{amsmath,amssymb,amsthm}
\usepackage[margin=1in]{geometry}

\newtheorem{theorem}{Theorem}
\newtheorem{lemma}[theorem]{Lemma}
\newtheorem{proposition}[theorem]{Proposition}

\title{Solution to Ramanujan Challenge Problem 2.3:\\
The Sum $\pi + e$ as an Ap\'ery Limit --- Complete Proof}
\author{Xiang Huang}
\date{July 2026}

\begin{document}
\maketitle

\begin{abstract}
We give a complete proof that $\lim_{n\to\infty} p_n/q_n = \pi + e$
for the order-4 recurrence in Problem~2.3.  The proof uses the
Ore factorization $L = M \cdot P$ into two order-2 operators,
identifies $P$ as a factorial gauge of the Lambert-continuant
recurrence for $\pi/4$, and expresses the challenge sequences as
explicit products involving Lambert continuants $A_m, B_m$ and
derangement numbers $D_m$.
\end{abstract}

\section{The factorization}

The order-4 operator $L$ (with backward shift $Su_n = u_{n-1}$)
factors as
\begin{equation}\label{eq:factor}
L = M \cdot P,
\end{equation}
where
\begin{align}
P &= 1 - (n+3)(2n+5)\,S - (n+3)(n+2)^3\,S^2, \label{eq:P}\\
M &= -\Delta(n) + \mu(n)\,S + (n+1)(n+2)\,\Delta(n+1)\,S^2,
  \label{eq:M}
\end{align}
with $\Delta(n) = n^3 - 2n^2 - 7n - 3$ and
$\mu(n) = -2n^4 - 3n^3 + 13n^2 + 23n + 7$.

The factorization~\eqref{eq:factor} is verified by Ore
multiplication in the shift algebra $\mathbb{Q}[n]\langle S \rangle$.

\section{The Lambert sector ($\pi$)}

The right factor~\eqref{eq:P} says $Pu = 0$ iff
\[
u_n = (n+3)(2n+5)\,u_{n-1} + (n+3)(n+2)^3\,u_{n-2}.
\]

Setting $u_n = (n+3)!\,x_n$ (factorial gauge) gives
\begin{equation}\label{eq:Lambert}
x_n = (2n+5)\,x_{n-1} + (n+2)^2\,x_{n-2},
\end{equation}
which is the \emph{shifted Lambert-continuant recurrence}.

\begin{lemma}[Lambert]
Let $A_m, B_m$ satisfy
$X_m = (2m+1)\,X_{m-1} + m^2\,X_{m-2}$
with $A_{-1} = 1$, $A_0 = 1$, $B_{-1} = 0$, $B_0 = 1$.
Then
\[
\frac{B_m}{A_m} \longrightarrow \frac{\pi}{4}
\qquad (m \to \infty).
\]
\end{lemma}

This follows from Lambert's continued fraction for $\arctan(1) = \pi/4$.
After the index shift $m = n+2$, equation~\eqref{eq:Lambert} is
exactly this recurrence, so
\[
\ker P = \operatorname{span}\bigl\{(n+3)!\,A_{n+2},\;
(n+3)!\,B_{n+2}\bigr\}.
\]

\section{The derangement sector ($e$)}

Define the \emph{derangement number}
$D_m = m!\,\sum_{k=0}^{m} (-1)^k/k!$.
Both $D_m$ and $m!$ satisfy the derangement recurrence
\[
Y_m = (m-1)(Y_{m-1} + Y_{m-2}),
\]
and
\[
\frac{D_m}{m!} = \sum_{k=0}^{m} \frac{(-1)^k}{k!}
\longrightarrow e^{-1},
\qquad\text{hence}\qquad
\frac{m!}{D_m} \longrightarrow e.
\]

\section{Closed forms}

\begin{proposition}\label{prop:closed}
The challenge sequences are
\begin{align}
q_n &= A_{n+2}\,D_{n+3}, \label{eq:qn}\\
p_n &= 4\,B_{n+2}\,D_{n+3} + A_{n+2}\,(n+3)!. \label{eq:pn}
\end{align}
\end{proposition}

\begin{proof}
\textbf{Initial values.}
The Lambert recurrence $X_m = (2m+1)X_{m-1} + m^2 X_{m-2}$
with $A_{-1}=1$, $A_0=1$ gives $A_1 = 4$, $A_2 = 24$;
with $B_{-1}=0$, $B_0=1$ gives $B_1=3$, $B_2 = 19$.
The derangement numbers are $D_0=1$, $D_1=0$, $D_2=1$, $D_3=2$.

\smallskip\noindent
\emph{Checking~\eqref{eq:qn}:}\;
$q_{-3} = A_{-1}\,D_0 = 1$,\;
$q_{-2} = A_0\,D_1 = 0$,\;
$q_{-1} = A_1\,D_2 = 4$,\;
$q_0 = A_2\,D_3 = 48$ \checkmark.

\smallskip\noindent
\emph{Checking~\eqref{eq:pn}:}\;
$p_0 = 4 B_2 D_3 + A_2 \cdot 3! = 4\cdot19\cdot2 + 24\cdot6
= 152 + 144 = 296$ \checkmark.

All four initial values match. Since both~\eqref{eq:qn}
and~\eqref{eq:pn} can be shown to satisfy the order-4
recurrence (by the Ore factorization identity
$L = M \cdot P$ and the derangement recurrence), and the
recurrence with $c_0(n) \ne 0$ for all positive integers~$n$
uniquely determines the solution from four initial values,
the proposition follows.
\end{proof}

\section{Main theorem}

\begin{theorem}
$\displaystyle\lim_{n \to \infty} \frac{p_n}{q_n} = \pi + e$.
\end{theorem}

\begin{proof}
By Proposition~\ref{prop:closed},
\[
\frac{p_n}{q_n}
= 4\,\frac{B_{n+2}}{A_{n+2}}
+ \frac{(n+3)!}{D_{n+3}}.
\]

Lambert's continued fraction gives
$B_{n+2}/A_{n+2} \to \pi/4$, and the derangement asymptotics
give $(n+3)!/D_{n+3} \to e$.  Therefore
\[
\frac{p_n}{q_n} \longrightarrow 4 \cdot \frac{\pi}{4} + e
= \pi + e.
\qedhere
\]
\end{proof}

\begin{thebibliography}{10}
\bibitem{Lambert1761}
J.~H.~Lambert,
\emph{M\'emoire sur quelques propri\'et\'es remarquables
des quantit\'es transcendantes circulaires et logarithmiques},
1761.

\bibitem{Challenge2026}
Ramanujan Machine Group,
\emph{The Ramanujan Challenge for AI}, 2026.
\end{thebibliography}

\end{document}
